%% CREATION D'UN VECTEUR TEMPS %%
%
\section{Cr\'eation d'un vecteur time � partir des vecteurs heure et minute}

\itemstyle{Op\'erateur~:} \operfmt{Cre\_time\_hm}
\index{Cre\_time\_hm}

%%%%%OBJET%%%%
\itemstyle{Objet~:}
Cette fonction cr\'ee un vecteur temps $Time$ (format \texttt{double},
secondes depuis 0h00 le jour de d�marrage de l'acquisition), � partir
de vecteurs heure et minute. Les secondes sont toutes suppos�es nulles.

%%%%%ARGUMENTS%%%%
\bigskip
%%%%%ENTREES%%%%
\begin{tabarg}
\multicolumn{3}{l}{\titretab {Entr\'ees}}\\
\hline
$H$ & Vecteur, {\tt int} & Heures\\
$M$ & Vecteur, {\tt int} & Minutes\\
\end{tabarg}

\medskip
%%%%%SORTIES%%%%
\begin{tabarg}
\multicolumn{3}{l}{\titretab {Sorties}}\\
\hline
$Time-out$ & Vecteur, {\tt double} & Vecteur temps en sortie\\
\end{tabarg}


%%%%%FORMULE%%%%
\itemstyle{Formule}


% calcul du vecteur time a partir de HH et MM
\[
time-out = \left( H * 3600. \right) + \left( M * 60. \right) 
\]


%%%%%COMMENTAIRE%%%%
\itemstyle{Commentaire~:}
Ne pas pas utiliser sur des s�ries passant
la limite de minuit, o� la valeur de l'heure d�croit
brusquement alors que le nombre de secondes depuis 0h00 le jour
du d�but de la s�rie (d�finition des "Time" dans AIDA) devrait
continuer � cro�tre.

%%%%%AUTEUR%%%%
\itemstyle{Auteur~:}
CNRM/GMEI/TRAMM

%%%%%REFERENCE%%%%
\itemstyle{R\'ef\'erences~:}

